% critique.tex — f165r, vv. 209–251 : commentaire critique (articles textuels tirés du commentaire de l'auteur ; répartition : repartition.md)
\cc{209}{ὦ τρισκακόδαιμον ὡς ἀπόλωλ᾽ Εὐριπίδης.}{La plupart des éditeurs modernes découpent cette réplique, attribuant la première partie à Euripide, puis Mnésiloque de lancer l'apostrophe : Εὐριπίδη κτλ. Cela les oblige à retoucher le texte pour faire de ὤ τρισκακοδαίμων le sujet de ἀπόλωλ(α) et changer Εὐριπίδης en Εὐριπίδη... 
Tout cela paraît bien tarabiscoté alors que le texte de notre unicodex se laisse comprendre facilement : « Ô triple malheureux ! ainsi périt Euripide », apostrophe qu'Euripide s'adresse à lui-même d'autant moins étonnante qu'il a déjà parlé de lui à la \nieme{3} personne (v. 77).
En outre, si l'attribution des répliques dans le manuscrit de Ravenne est souvent folklorique, la découpe des répliques, elle, est bien plus sûre. Philologues, réfrénez vos pulsions à vouloir tailler ainsi le vers.}
\cc{217}{ἢ μὴ διδόναι γ' αὐτὸν ὤφελόν ποτε.}{Le vers pèche contre le mètre.
Willems (voir sa note 4 en annexe), en guerre perpétuelle contre les retoucheurs de texte, rendait justice à ce qu'il croyait être le texte original. (En fait légèrement modifié : γ' ἑμαυτὸν à la place de γ' αὐτὸν. Willems a-t-il été trompé par Dindorf ?) 
Il s'en justifiait par des exemples qui ne laissent nul doute de l'epsilonesque correction qu'il aurait lui-même proposée s'il avait eu le texte des manuscrits sous les yeux : substituer γ' ἑαυτὸν à γ' αὐτὸν.\\ 
Au reste, comme l'avait souligné Willems, « διδόναι ἑαυτόν τινι est la formule ordinaire, et qui se lit partout : \italique{Cavaliers}, 740 ; Sophocle, \italique{Philoct}., 84 ; \italique{Trach}., 1117; Hérodote, VI, 108; Thucydide, II, 68; Démosthène, \italique{Sur la couronne}, 503d.~»}
\cc{230}{ἔχ᾽ ἀτρέμα σαυτὸν κἀνάκυπτε·}{C'est la leçon du \italique{Ravennas}. 
M. van Leeuwen, d'après Dobree, corrige : ἔχ᾽ ἀτρέμας αὐτοῦ, \italique{quia ferri nequit accusativus} ; lequel αὐτοῦ était à sa place en \italique{Oiseaux}, 1200, mais pas du tout ici. 
Hérodote a pourtant dit, en parlant de l'armée athénienne : εἶχον ἀτρέμας σφέας αὐτοὺς ἵνα ἐτάχθησαν\nc{Hérodote, IX, 54 (aussi III, 79). — De même Platon, \italique{Epist.}, 348 B, écrit : ἐκποδὼν ἑαυτὸν ἔσχεν et Euripide, \italique{Iph. Taur.}, 1129 : ἐκποδὼν τοῦδ' ἔχειν μιάσματος.}. Hérodote le pouvait, mais non Aristophane, qui ailleurs s'est borné à ἔχ' ἀτρέμα\nc{Il est vrai qu'une des quatre fois les bons manuscrits portent ἔχ' ἀτρεμί (\italique{Nuées}, 261). Était-ce assez dégoûtant ? Heureusement que M. van Leeuwen y a mis bon ordre.}. Quoi ! voilà un comique qui, à quatre reprises (car on s'est donné la peine de compter), a écrit : \italique{reste coi}, et la cinquième fois il se permettrait d'écrire : \italique{tiens-toi coi}. 
À quoi servirait la statistique si l'on tolérait de pareilles licences ?
(Note vengeresse n°5 de Willems.)}
\cc{234}{βούλει θεάσασθαι σεαυτόν; — εἰ δοκεῖ, φέρε.}{Le vers péchant contre le mètre, la correction la plus convaincante me semble être celle de Fritzsche qui substitue αὑτόν au σεαυτόν, faisant l'hypothèse que σεαυτόν aura été suggéré à un scribe inattentif par la proximité avec ὁρᾷς σεαυτόν du vers suivant.}
